Pozície nájdených primerov boli uložené do poľa v tvare:
\begin{verbatim}
[[], [65], [55], [], [64], [], [55], [74, 73], [64], [75], ...]
\end{verbatim}

Každý prvok poľa zodpovedá jednému čítaniu a obsahuje pozície, na ktorých sa primer nachádza v danom čítaní. Z výsledkov vyplýva, že v nie každom reťazci bol primer identifikovaný, zatiaľ čo v niektorých čítaniach sa vyskytol viackrát.
Tento jav môže byť spôsobený tým, že niektoré čítania obsahujú komplementárnu sekvenciu alebo boli prečítané v opačnom smere. Ďalším dôvodom môže byť pretrhnutie DNA reťazca v rôznych miestach alebo príliš nízka podobnosť primeru vzhľadom na nastavený prah zarovnania.
